% ============ 第1章 §1.4 插值理论 (原书 pp.14-19) ============
\section{插值理论}\label{sec:1.4}
\index{M@Marcinkiewicz 插值定理}




解析学为何复杂, 理由之一是调和级数发散. 也就是说,
\[
\frac{1}{1} + \frac{1}{2} + \cdots + \frac{1}{n} + \cdots = \infty,
\]
所以各种问题就变得困难. 关于这个级数为何复杂, 一个说明是: 若 $1 < p < \infty$,
\[
\frac{1}{1^{p}} + \frac{1}{2^{p}} + \cdots + \frac{1}{n^{p}} + \cdots < \infty.
\]
也就是说, $p = 1$ 的情形是特别的. 更进一步说, 可以认为 $p = 1$ 的情形是病态的. 然而, 在解析学中, 与
\[
\frac{1}{1} + \frac{1}{2} + \cdots + \frac{1}{n} + \cdots = \infty
\]
类似的现象, 例如
\[
\int_{1}^{\infty} \frac{\dd t}{t} = [\log t]_{1}^{\infty} = \infty
\]
会出现, 所以调和级数发散这一事实是重要的.

来稍微定量地把握调和级数发散这一问题. 具体地, 用数学归纳法可以证明
\[
\frac{1}{1} + \frac{1}{2} + \cdots + \frac{1}{n} > \log(n+1) \quad (n \in \N).
\]
这里重要的是, 调和级数虽然发散, 但其发散程度非常弱. 也就是说, 第 $n$ 部分和大于 $\log n$, 但并不是大于 $n$ 的任意幂 $cn^{a}$, $a, c > 0$, 这是应该注意的事实. 也就是说, 虽然
\[
\frac{1}{1} + \frac{1}{2} + \cdots + \frac{1}{n} + \cdots = \infty,
\]
但这个无穷大在某种意义上是很弱的发散, 或许应该称之为``准有限''. 这里, 我们说明: 当有两个``准有限''时, 能够产生有限的插值理论. 下面给出的定理是被称为 Marcinkiewicz 插值定理的 1930 年代的定理.

设 $S$ 为从 $\R^{N}$ 到 $\R^{N}$ 的实线性映射. 将 $a \in \R^{N}$ 经 $S$ 的像记为 $Sa = \{(Sa)_{n}\}_{n=1}^{N}$. 对 $t > 0$, 假定
\[
\# \{1 \leq n \leq N : |(Sa)_{n}| > t\} \leq \min(t^{-1}\|a\|_{\ell^{1}}, t^{-q}\|a\|_{\ell^{q}}^{q})
\]
成立. 来考察
\[
\|Sa\|_{\ell^{p}}^{p} = p\int_{0}^{\infty} t^{p-1} \# \{1 \leq n \leq N : |(Sa)_{n}| > t\} \dd t.
\]
对 $a \in \R^{N}$ 与 $t > 0$, 将 $c^{t} = \{c_{k}^{t}\}_{k=1}^{N}$ 定义为: 在 $a$ 的分量中, 绝对值大于 $t$ 的保持原样, 绝对值不超过 $t$ 的替换为 $0$. 进一步, 令 $b^{t} = a - c^{t}$. 此时, 由于 $(Sa)_{n} = (Sb^{t})_{n} + (Sc^{t})_{n}$, 由三角不等式,
\begin{align*}
\{1 \leq n \leq N : |(Sb^{t})_{n}| \leq 2^{-1}t\} \cap \{1 \leq n \leq N : |(Sc^{t})_{n}| \leq 2^{-1}t\} \subset \{1 \leq n \leq N : |(Sa)_{n}| \leq t\}.
\end{align*}
因此, 设 $1 < p < q$, 由假设得
\begin{align*}
& \# \{1 \leq n \leq N : |(Sa)_{n}| > t\} \\
& \leq \# \{1 \leq n \leq N : |(Sb^{t})_{n}| > 2^{-1}t\} + \# \{1 \leq n \leq N : |(Sc^{t})_{n}| > 2^{-1}t\} \\
& \leq 2^{q}t^{-q}\|b^{t}\|_{\ell^{q}}^{q} + 2t^{-1}\|c^{t}\|_{\ell^{1}}.
\end{align*}
这里,
\begin{align*}
\int_{0}^{\infty} t^{p-2}\|c^{t}\|_{\ell^{1}} \dd t & = \int_{0}^{\infty} t^{p-2}\sum_{j=1}^{N} |a_{j}|\chi_{(t,\infty)}(|a_{j}|) \dd t \\
& = \sum_{j=1}^{N} |a_{j}|\int_{0}^{\infty} t^{p-2}\chi_{(t,\infty)}(|a_{j}|) \dd t \\
& = \frac{1}{p-1}\sum_{j=1}^{N} |a_{j}|^{p}
\end{align*}
以及
\begin{align*}
\int_{0}^{\infty} t^{p-q-1}\|b^{t}\|_{\ell^{q}}^{q} \dd t & = \int_{0}^{\infty} t^{p-q-1}\sum_{j=1}^{N} |a_{j}|^{q}\chi_{[0,t]}(|a_{j}|) \dd t \\
& = \sum_{j=1}^{N} |a_{j}|^{q}\int_{0}^{\infty} t^{p-q-1}\chi_{[0,t]}(|a_{j}|) \dd t \\
& = \sum_{j=1}^{N} |a_{j}|^{q}\int_{|a_{j}|}^{\infty} t^{p-q-1} \dd t \\
& = \frac{1}{q-p}\sum_{j=1}^{N} |a_{j}|^{p}
\end{align*}
所以,
\begin{equation}\tag{1.8}\label{eq:1.8}
\|Sa\|_{\ell^{p}} \leq \left(\frac{2p}{p-1} + \frac{2^{q}p}{q-p}\right)^{\frac{1}{p}}\|a\|_{\ell^{p}}
\end{equation}
成立. 由于 $a$ 是任意的,
\[
\|S\|_{\ell^{p} \to \ell^{p}} \leq \left(\frac{2p}{p-1} + \frac{2^{q}p}{q-p}\right)^{\frac{1}{p}}.
\]
这样, 从关于 $\ell^{1}$, $\ell^{q}$ 已知的不等式得到 $\ell^{p}$, $1 < p < q$ 的不等式的理论, 称为插值理论.

当 $q = \infty$ 时, 将假设改写为 $\|Sa\|_{\ell^{\infty}} \leq 2^{-1}\|a\|_{\ell^{\infty}}$ 以及
\[
\# \{1 \leq n \leq N : |(Sa)_{n}| > t\} \leq t^{-1}\|a\|_{\ell^{1}}.
\]
此时要使用
\begin{align*}
\{1 \leq n \leq N : |(Sb^{t})_{n}| > 2^{-1}t\} & = \{1 \leq n \leq N : 2^{-1}t \geq |(Sb^{t})_{n}| > 2^{-1}t\} \\
& = \emptyset.
\end{align*}
无论 $q = \infty$ 还是 $q \neq \infty$, 都可以在
\[
\# \{1 \leq n \leq N : |(Sa)_{n}| > t\} \leq \min(t^{-1}\|a\|_{\ell^{1}}, t^{-q}\|a\|_{\ell^{q}}^{q})
\]
或 $\|Sa\|_{\ell^{\infty}} \leq \|a\|_{\ell^{\infty}}$ 以及
\[
\# \{1 \leq n \leq N : |(Sa)_{n}| > t\} \leq t^{-1}\|a\|_{\ell^{1}}
\]
的位置, 改用对适当的常数 $K > 0$ 的不等式
\[
\# \{1 \leq n \leq N : |(Sa)_{n}| > t\} \leq K\min(t^{-1}\|a\|_{\ell^{1}}, t^{-q}\|a\|_{\ell^{q}}^{q})
\]
或 $\|Sa\|_{\ell^{\infty}} \leq K\|a\|_{\ell^{\infty}}$ 以及
\[
\# \{1 \leq n \leq N : |(Sa)_{n}| > t\} \leq Kt^{-1}\|a\|_{\ell^{1}}.
\]
此时, 根据 $q$ 是有限还是无限, 使用 $K^{-1}S$ 或 $(2K)^{-1}S$ 来使用插值定理.

\begin{theorem}[Marcinkiewicz 插值定理]\label{thm:1.5}
(1) 设 $1 < p < q < \infty$. 对任意 $t > 0$ 与数列 $\{a_{n}\}_{n=1}^{N}$, 不等式
\[
\# \{1 \leq n \leq N : |(Sa)_{n}| > t\} \leq \min(Kt^{-1}\|a\|_{\ell^{1}}, K^{-q}t^{-q}\|a\|_{\ell^{q}}^{q})
\]
成立的线性算子 $S: \R^{N} \to \R^{N}$, 存在仅依赖于 $p$, $q$ 的适当常数 $c_{p,q} > 0$, 使得 $\|S\|_{\ell^{p} \to \ell^{p}} \leq c_{p,q}K$ 成立.

(2) 设 $1 < p < \infty$. 对满足不等式 $\|S\|_{\ell^{\infty} \to \ell^{\infty}} \leq K$ 且对任意数列 $\{a_{n}\}_{n=1}^{N}$,
\[
\# \{1 \leq n \leq N : |(Sa)_{n}| > t\} \leq Kt^{-1}\|a\|_{\ell^{1}} \quad (t > 0)
\]
成立的线性算子 $S: \R^{N} \to \R^{N}$, 存在仅依赖于 $p$ 的适当常数 $c_{p} > 0$, 使得 $\|S\|_{\ell^{p} \to \ell^{p}} \leq c_{p}K$ 成立.
\end{theorem}

来举出满足定理假设的算子的具体例子.

\begin{example}\label{exa:1.8}
(1) 设数列 $a = \{a_{n}\}_{n=1}^{N}$ 已给出. 对满足 $|a_{n}| > t$ 的 $n = 1, 2, \ldots, N$, 有 $t^{-p}|a_{n}|^{p} > 1$. 因此,
\[
\# \{1 \leq n \leq N : |a_{n}| > t\} = \sum_{n=1}^{N} \chi_{(t,\infty)}(|a_{n}|) \leq t^{-p}\|a\|_{\ell^{p}}^{p}
\]
成立.

(2) 设 $1 \leq p \leq \infty$. 定义域与值域都考虑 $\ell^{p}$-范数, 当线性算子 $S: \R^{N} \to \R^{N}$ 的范数不超过 $1$ 时, 对所有数列 $a$, $\|Sa\|_{\ell^{p}} \leq \|a\|_{\ell^{p}}$ 成立, 从而 $\# \{1 \leq n \leq N : |(Sa)_{n}| > t\} \leq t^{-p}\|a\|_{\ell^{p}}^{p}$ 成立. 特别地, 若对 $p = 1$, $q$ 双方这个条件都满足, 则 Marcinkiewicz 插值定理的假设成立.
\end{example}

\begin{example}\label{exa:1.9}
对实数列 $\{a_{n}\}_{n=1}^{2N}$, 令
\[
b_{n} = \frac{a_{1} + a_{2} + \cdots + a_{2n}}{2n}, \quad (n = 1, 2, \ldots, N), \quad b_{N+1} = \cdots = b_{2N} = 0.
\]
考虑对应 $T_{N}: \{a_{n}\}_{n=1}^{2N} \mapsto \{b_{n}\}_{n=1}^{2N}$.

(1) 对 $n = 1, 2, \ldots, 2N$, $|b_{n}| \leq \frac{1}{2n}\|a\|_{\ell^{1}}$ 成立, 所以
\[
\# \left\{k \in \{1, 2, \ldots, 2N\} : |b_{k}| > \frac{1}{2n}\|a\|_{\ell^{1}}\right\} \leq \# \{1, 2, \ldots, n-1\} = n-1
\]
成立.

(2) 若 $\lambda \geq \frac{1}{2}$, 则出现在左边的集合是空集, 所以
\[
\# \left\{k \in \{1, 2, \ldots, 2N\} : |b_{k}| > \lambda\|a\|_{\ell^{1}}\right\} = \#\emptyset = 0.
\]

(3) 另外, 由于出现在左边的集合是 $\{1, 2, \ldots, 2N\}$ 的子集, 若 $\lambda \leq \frac{1}{2N}$,
\[
\# \left\{k \in \{1, 2, \ldots, 2N\} : |b_{k}| > \lambda\|a\|_{\ell^{1}}\right\} \leq 2N \leq \frac{1}{\lambda}.
\]

(4) 若 $\frac{1}{2N} < \lambda < \frac{1}{2}$, 则存在 $n = 1, 2, \ldots, N-1$ 使得 $\frac{1}{2(n+1)} \leq \lambda < \frac{1}{2n}$, 所以
\begin{align*}
\# \left\{k \in \{1, 2, \ldots, 2N\} : |b_{k}| > \lambda\|a\|_{\ell^{1}}\right\} & \leq \# \left\{k \in \{1, 2, \ldots, 2N\} : |b_{k}| > \frac{1}{2n+2}\|a\|_{\ell^{1}}\right\} \\
& = \# \{1, 2, \ldots, n\} = n \leq 2n \leq \frac{1}{\lambda}.
\end{align*}
综上所述, 对任意 $\lambda > 0$,
\[
\# \left\{k \in \{1, 2, \ldots, 2N\} : |b_{k}| > \lambda\|a\|_{\ell^{1}}\right\} \leq \frac{1}{\lambda}
\]
成立. 也就是说, $T_{N}$ 满足定理~\ref{thm:1.5} 的假设.
\end{example}

来实际说明用定理能得到什么.

\begin{example}\label{exa:1.10}
设 $1 \leq p \leq \infty$. 在定义域、值域 $\R^{2N}$ 上配备 $\ell^{p}$-范数, 考虑例~\ref{exa:1.9} 中考虑的算子 $T_{N}$.

(1) 设 $p = 1$. 对实数列 $a = \{a_{n}\}_{n=1}^{2N}$,
\begin{align*}
\|T_{N}a\|_{\ell^{1}} & = \sum_{n=1}^{N} \frac{|a_{1} + a_{2} + \cdots + a_{2n}|}{2n} \\
& \leq \sum_{n=1}^{N} \frac{|a_{1}| + |a_{2}| + \cdots + |a_{2n}|}{2n} \\
& \leq \sum_{n=1}^{N} \frac{|a_{1}| + |a_{2}| + \cdots + |a_{2N}|}{2n} \\
& = \sum_{n=1}^{N} \frac{\|a\|_{\ell^{1}}}{2n}
\end{align*}
由于存在 $a = e_{1}$ 使上式等号成立, 所以 $\|T_{N}\|_{\ell^{1} \to \ell^{1}} = \sum_{n=1}^{N} \frac{1}{2n}$.

(2) 设 $p = \infty$. 平均不会增大 $\ell^{\infty}$-范数, 所以 $\|T_{N}\|_{\ell^{\infty} \to \ell^{\infty}} = 1$ 成立.

(3) 设 $1 < p < \infty$. 由例~\ref{exa:1.9} 与 Marcinkiewicz 插值定理, 存在不依赖于 $N$ 的常数 $c_{p}$, 使得 $\|T_{N}\|_{\ell^{p} \to \ell^{p}} \leq c_{p}$ 成立.

(4) 实际计算可知, $\|T_{1}\|_{\ell^{2} \to \ell^{2}} = \frac{1}{\sqrt{2}}$, $\|T_{2}\|_{\ell^{2} \to \ell^{2}} = \frac{1 + \sqrt{5}}{4}$ (参照例~\ref{exa:1.5}), $\|T_{3}\|_{\ell^{2} \to \ell^{2}} = 0.86\cdots$. 这样, 对 $j = 1, 2, \ldots$, $\|T_{j}\|_{\ell^{2} \to \ell^{2}}$ 递增, 但可以看出不会超过一定值.
\end{example}

\begin{exercise}\label{exer:1.17}
对实数列 $\{a_{n}\}_{n=1}^{N}$, 令
\[
b_{n} = \frac{a_{1} + a_{2} + \cdots + a_{n}}{n}, \quad (n = 1, 2, \ldots, N),
\]
考虑数列 $\{b_{n}\}_{n=1}^{N}$. 设 $1 \leq p \leq \infty$. 在定义域、值域 $\R^{N}$ 上配备 $\ell^{p}$-范数, 考虑这个对应 $U_{N}: \{a_{n}\}_{n=1}^{N} \mapsto \{b_{n}\}_{n=1}^{N}$.

(1) 证明 $\|U_{N}\|_{\ell^{1} \to \ell^{1}} = \sum_{n=1}^{N} \frac{1}{n}$.

(2) 证明 $\|U_{N}\|_{\ell^{\infty} \to \ell^{\infty}} = 1$.

(3) 设 $\lambda > 0$. 证明 $\# \{k \in \{1, 2, \ldots, N\} : |b_k| > \lambda\|a\|_{\ell^{1}}\} \leq \frac{1}{\lambda}$.

(4) 设 $1 < p < \infty$. 证明存在不依赖于 $N$ 的常数 $c_{p} > 0$, 使得 $\|U_{N}\|_{\ell^{p} \to \ell^{p}} \leq c_{p}$ 成立.

(5) 实际计算并比较 $\|U_{1}\|_{\ell^{2} \to \ell^{2}}$, $\|U_{2}\|_{\ell^{2} \to \ell^{2}}$, $\|U_{3}\|_{\ell^{2} \to \ell^{2}}$. 其中, 关于 $\|U_{3}\|_{\ell^{2} \to \ell^{2}}$, 只要导出能看出这三者大小关系程度的信息即可.
\end{exercise}
